% Generated draft: requires mathematical and source review.
\documentclass[11pt]{article}
\usepackage[margin=25mm]{geometry}
\usepackage{fontspec}
\setmainfont{Latin Modern Roman}
\usepackage{amsmath,amssymb,mathtools,unicode-math}
\setmathfont{Latin Modern Math}
\usepackage{enumitem,needspace,xurl,hyperref}
\setlength{\parindent}{0pt}
\setlength{\parskip}{5pt}
\setlength{\emergencystretch}{3em}
\newcommand{\N}{\mathbb{N}}
\newcommand{\Z}{\mathbb{Z}}
\newcommand{\Q}{\mathbb{Q}}
\newcommand{\R}{\mathbb{R}}
\newcommand{\C}{\mathbb{C}}
\newcommand{\F}{\mathbb{F}}
\newcommand{\A}{\mathbb{A}}
\newcommand{\PP}{\mathbb{P}}
\newcommand{\E}{\mathbb{E}}
\newcommand{\eps}{\varepsilon}
\newcommand{\dd}{\,\mathrm{d}}
\begin{document}
\section*{UnsolvedMath ID: 51}
\textbf{Source problem code:} catalogue entry
\subsection{Definitions and mathematical statement}
Let $S^2=\{x\in\R^3:\|x\|=1\}$. For an integer $n\ge1$, let
\[
 \mathcal C_n=\{(x_1,\ldots,x_n)\in(S^2)^n:x_i\ne x_j\text{ for }i\ne j\}
\]
be the space of configurations of $n$ distinct points. In the usual pairwise chordal-distance convention, a specified potential $\Phi:(0,2]\to\R$ gives the energy
\[
 E_\Phi(x_1,\ldots,x_n)=\sum_{1\le i<j\le n}\Phi(\|x_i-x_j\|),\qquad
 e_\Phi(n)=\inf_{\mathbf x\in\mathcal C_n}E_\Phi(\mathbf x).
\]
For each intended potential, the exact minimization request would be to determine $e_\Phi(n)$ and configurations $\mathbf x$ satisfying
\[
 \forall\mathbf y\in\mathcal C_n,\qquad
 E_\Phi(\mathbf x)\le E_\Phi(\mathbf y),
\]
for general $n$, with configurations optionally identified under permutations and orthogonal transformations of $\R^3$. The source's question asks for optimal point arrangements for ``various potential functions,'' but does not specify which potentials, whether the energy is pairwise, whether chordal or geodesic distance is used, whether coincidences are allowed, or whether exact optimization or an algorithmic approximation is intended. The display is therefore a concrete conventional framework, not a uniquely determined repair of the source. With no regularity or repulsiveness assumptions on $\Phi$, attainment of the infimum on the distinct-point configuration space must not be assumed.

The supplied background includes the Thomson problem among the examples. In the standard Coulomb interpretation that corresponds to $\Phi(r)=1/r$ in the displayed framework, but the record does not limit its target to that potential. The background also mentions known small-$n$ and highly symmetric examples and connections to crystallography and coding theory.

The literature triage checked 17 August 2026 says that universal-optimality results cover certain highly symmetric configurations and classes of potentials, without specifying those configurations or classes. Such a result concerns simultaneous minimization over its particular class; it does not supply one arrangement minimizing every arbitrary potential. The triage labels the row partially solved and calls its broad wording a programme rather than a single proposition, with no universal solution covering it. The registry's main status is open. The record attributes the problem to Stephen Smale in 1998 and labels it his seventh problem, but supplies only a secondary overview, not a precise original formulation or computational accuracy requirement.
\subsection{Short English statement}
For a specified interaction potential, determine the configurations of $n$ points on the sphere with globally minimal energy. The record ranges over unspecified potentials and leaves the metric and exact computational target open, so it describes a programme rather than one precise optimization problem.
\subsection{Sources}
\begin{itemize}
\item Local UnsolvedMath registry, permanent ID $51$, source code \texttt{catalogue entry}, ``Smale's 7th Problem: Distribution of Points on the 2-Sphere,'' attributed there to Stephen Smale, 1998; including the supplied literature triage checked 17 August 2026.
\item \emph{Smale's Problems}, MathWorld, accessed 17 August 2026 according to the record, \url{https://mathworld.wolfram.com/SmalesProblems.html}. This is the supplied secondary reference; no original problem citation or precise universal-optimality reference is supplied.
\end{itemize}
\end{document}
